%%%%%%%%%%%%%%%%%%%%%%%%%%%%%%%%%%%%%%%%%%%%%%%%%%%%%%%%%%%%%%%%%%%%%%%%%%%%%%%
%% Apêndice A --- Registro de verificação do hardware
%%
%% Absorve o registro de engenharia do relatório de verificação de hardware:
%% correções no esquemático, defeitos de geração corrigidos e abertos, versões.
%% O material de PROJETO (cadeia de geração, placa, PDN) fica no Capítulo 3.
%%
%% PENDENTE DE TRANSCRIÇÃO: as tabelas completas de defeitos (itens 1--21) e de
%% comparação entre versões v1/v2 estão em docs/relatorio/verificacao-hardware.tex
%% e devem ser transportadas na íntegra. A estrutura abaixo é a do destino.
%%%%%%%%%%%%%%%%%%%%%%%%%%%%%%%%%%%%%%%%%%%%%%%%%%%%%%%%%%%%%%%%%%%%%%%%%%%%%%%

\chapter{Registro de verificação do hardware}
\label{ape:hardware}

Este apêndice reúne o registro de engenharia da cadeia de geração descrita na
Seção~\ref{sec:dev-hardware}. Ele é mantido no trabalho por uma razão de método: a
verificação só é informativa se os defeitos que ela encontrou forem publicados junto com os
que ela não encontrou.

Nenhum resultado aqui provém de bancada. Todos os valores são de geração, verificação
geométrica ou simulação, e o registro não reivindica validação experimental. Cada afirmação
tem associada a grandeza que a sustenta; itens sem medição associada aparecem como
\textsf{não verificado}.

\section{Correções no esquemático e defeitos corrigidos}
\label{sec:ape-corrigidos}

Os defeitos identificados e corrigidos durante a geração constam do relatório de
verificação com o valor antigo, o valor corrigido e a causa. Incluem, entre outros, erro de
referencial angular de \emph{pad} que produzia \num{42} curtos no DRC, e inversão de sinal
de rotação de \emph{footprint}.

\section{Defeitos abertos}
\label{sec:ape-abertos}

\emandamento{Os defeitos abaixo permanecem abertos e restringem afirmações específicas.}

\begin{itemize}
  \item \textbf{A placa não parafusa nos bosses do invólucro.} Na v1, o furo M3 exigiria
        material até \SI{28,2}{\milli\meter} na diagonal e o canto termina em
        \SI{27,5}{\milli\meter}; na v2, os recortes coincidem com os eixos dos bosses. A
        fixação passa a ser o assento pelo aro mais três furos M2 próprios.
  \item \textbf{Corredor do cabo.} Com \SI{6,0}{\milli\meter}, a v1 perde a saída de terra
        de dois componentes, com três itens soltos no DRC.
  \item \textbf{Altura do conjunto definida pelos conectores.} Decidir se o cabeçalho de
        gravação permanece populado no produto.
  \item \textbf{Divisor do NTC satura no frio}, atingindo \SI{3,22}{\volt} a
        \SI{-40}{\celsius}, acima da faixa útil do conversor. Abaixo de aproximadamente
        \SI{-35}{\celsius} os limiares de curto e de circuito aberto não disparam.
  \item \textbf{Numeração de pino do regulador varia por fabricante} ---
        \textsf{não verificado}; conferir contra a folha de dados antes de fabricar.
  \item \textbf{Furo de conector no invólucro sem conector correspondente} na netlist.
  \item \textbf{Rabicho de antena ausente da lista de materiais}; a ligação do conector
        do módulo de rádio ao painel não está especificada.
  \item \textbf{Nenhum código lê a tensão da bateria.} O caminho de hardware existe; os
        limiares do requisito de segurança correspondente não estão implementados.
\end{itemize}

\section{Limitações declaradas das análises}
\label{sec:ape-limitacoes}

A análise modal assume ASA isotrópico e maciço, hipótese otimista: a peça impressa com
preenchimento parcial é anisotrópica e menos rígida. O amortecimento do material não foi
medido, e sim herdado de análise anterior do projeto. As demais limitações constam da
Seção~\ref{sec:conc-alcance}.
